\section{Results}

\begin{table}[t]
\centering
\footnotesize
\begin{tabular}{p{0.24\linewidth}p{0.18\linewidth}p{0.22\linewidth}p{0.24\linewidth}}
\toprule
Condition & Utility & Exposure / manifestation & Scope caveat \\
\midrule
No defense & Clean utility 0.3733; attacked utility 0.3333 & Exposed poisoned retrieval 0.4667; attack manifestation 0.2533 & Comparator is adapted, not an official AgentPoison reproduction \\
FlowFence-Lite & Clean utility 0.36; attacked utility 0.3867 & Raw poisoned retrieval 0.44 internally; exposed poisoned retrieval and manifestation 0.0 & Quarantine plus canonical action writeback on saved P0 setting \\
Static keyword filter & Known-trigger weak comparator & Exposed poisoned retrieval and manifestation 0.0 on same-axis trigger & Strong caveat baseline; does not establish robustness to paraphrased or indirect attacks \\
Held-out instruction stress & FlowFence attacked utility 0.2533; static filter 0.1733 & Static keyword attack manifestation 0.24; FlowFence manifestation 0.0 & Same retrieval anchor only; supports blocklist-brittleness caveat, not broad held-out robustness \\
\bottomrule
\end{tabular}
\caption{P0 retrieval-memory containment evidence. The table summarizes committed high-level P0 artifacts and should be read as adapted AgentPoison comparator evidence only, not official AgentPoison reproduction.}
\label{tab:p0_agentpoison}
\end{table}

\begin{table*}[t]
\centering
\footnotesize
\begin{tabular}{p{0.16\linewidth}p{0.1\linewidth}p{0.1\linewidth}p{0.09\linewidth}p{0.09\linewidth}p{0.36\linewidth}}
\toprule
Group & Runs & Success & Raw leak & Ext. leak & Main result \\
\midrule
Canonical coverage & 252/252 & 0.964 & 2.829 & 0.325 & MiniMax-backed multi-agent synthetic-runtime coverage over 3 topologies, 7 attacks, 4 defenses, and 3 seeds; cascade mean 4.179; privilege reach mean 2.679 \\
FlowFence subset & 63/63 clean & 1.0 & 0.0 & 0.0 & Clean across seeds 1, 2, and 3; cascade can be non-zero because quarantine is still logged as containment activity \\
FlowFence vs no defense & 63 comparisons & 7 improves / 56 ties & 42 improves / 21 ties & 32 improves / 31 ties & No underperformance reported in configured comparisons \\
FlowFence vs static ACL & 63 comparisons & 63 ties & 42 improves / 21 ties & 30 improves / 33 ties & Static ACL retains policy gaps \\
FlowFence vs prompt filter & 63 comparisons & 2 improves / 61 ties & 18 improves / 45 ties & 13 improves / 50 ties & Prompt filter remains vulnerable in configured indirect settings; postfix smoke is legacy/superseded \\
Topology effect & Coverage summary & -- & -- & -- & Observed across chain, star, and blackboard topologies \\
\bottomrule
\end{tabular}
\caption{Canonical MiniMax-only P1 MAS coverage over the synthetic deterministic runtime. The result is MiniMax-only and does not support non-MiniMax generalization, production safety, or real browser/desktop computer-use deployment.}
\label{tab:minimax_3seed}
\end{table*}

\begin{table*}[t]
\centering
\footnotesize
\begin{tabular}{p{0.18\linewidth}p{0.1\linewidth}p{0.1\linewidth}p{0.09\linewidth}p{0.09\linewidth}p{0.09\linewidth}p{0.27\linewidth}}
\toprule
Setting & Runs & Calls & Success & Raw leak & Ext. leak & Comparison summary \\
\midrule
Deterministic non-oracle & 540/540/0 & no & 1.0 & 0.0 & 0.0 & Zero oracle-annotation violations; held-out paraphrases included \\
Targeted MiniMax non-oracle & 72/72/0 & yes & 1.0 & 0.0 & 0.0 & Targeted MiniMax-backed synthetic-runtime validation; zero oracle-annotation violations \\
Deterministic vs no defense & 90 comps. & no & -- & 81 improves / 9 ties & 81 improves / 9 ties & No underperformance on configured held-out comparisons \\
Deterministic vs static ACL & 90 comps. & no & -- & 66 improves / 24 ties & 48 improves / 42 ties & Static ACL policy gaps remain \\
Deterministic vs prompt filter & 90 comps. & no & -- & 54 improves / 36 ties & 54 improves / 36 ties & Prompt-filter paraphrase failures recorded \\
Targeted MiniMax vs no defense & 18 comps. & yes & -- & 15 improves / 3 ties & 10 improves / 8 ties & No underperformance on configured held-out comparisons \\
Targeted MiniMax vs static ACL & 18 comps. & yes & -- & 15 improves / 3 ties & 9 improves / 9 ties & Static ACL retains held-out policy gaps \\
Targeted MiniMax vs prompt filter & 18 comps. & yes & -- & 16 improves / 2 ties & 10 improves / 8 ties & MiniMax-only non-oracle held-out validation over the synthetic deterministic MAS runtime \\
No-semantic-pattern ablation & Det. subset & no & tied & higher than non-oracle & tied & Semantic patterns contribute to raw-leakage containment; policy/fanout/safe-view still matter \\
\bottomrule
\end{tabular}
\caption{Non-oracle held-out validation and mechanism ablation. The targeted run is MiniMax-only and synthetic-runtime; the evidence does not prove arbitrary attack robustness, non-MiniMax generalization, production safety, or real computer-use behavior.}
\label{tab:nonoracle}
\end{table*}


\paragraph{P0 retrieval-memory containment.}
Table~\ref{tab:p0_agentpoison} shows that no defense creates poisoned-retrieval exposure and non-zero attack manifestation on the adapted AgentPoison comparator. FlowFence-Lite quarantine with canonical action writeback reduces exposed poisoned retrieval and attack manifestation to zero on the saved P0 axis. This should be read as adapted retrieval-memory containment evidence, not official AgentPoison reproduction. The static keyword row is an important caveat: a known-trigger same-axis filter is strong in this setting.

\paragraph{Deterministic synthetic evidence.}
The strengthened deterministic MAS benchmark records topology-dependent propagation and shows FlowFence-Lite improving over prompt filtering on indirect synthetic attacks. This supports the benchmark mechanism but remains deterministic synthetic evidence.

\paragraph{MiniMax-backed coverage.}
Table~\ref{tab:minimax_3seed} summarizes the canonical MiniMax-only P1 MAS coverage over the synthetic deterministic runtime. The matrix completed 252/252 configured runs after retrying one transient timeout. The FlowFence subset is clean across 63/63 configured FlowFence runs, with task success 1.0 and raw/external leakage 0.0. FlowFence improves or ties no defense, static ACL, and prompt filtering on configured raw/external leakage comparisons. The coverage records \texttt{topology\_effect\_observed=true}.

\paragraph{Seed stability.}
Across seeds 1, 2, and 3, each seed has 21/21 clean FlowFence runs. This supports seed-level stability for the configured FlowFence subset, not broader robustness.

\paragraph{Non-oracle held-out validation.}
Table~\ref{tab:nonoracle} addresses the oracle-annotation concern. The deterministic non-oracle validation completed 540/540 runs with no failures; \texttt{flowfence\_lite\_nonoracle} had task success 1.0, raw leakage 0.0, external leakage 0.0, and zero oracle-annotation violations. The targeted MiniMax-backed synthetic-runtime validation completed 72/72 runs with the same clean non-oracle FlowFence subset. The non-oracle variant improves or ties no defense, static ACL, and prompt filtering on configured held-out raw/external leakage comparisons.

\paragraph{Mechanism ablation.}
The no-semantic-pattern ablation tied the full non-oracle variant on external leakage and task success but had higher raw leakage in deterministic results. This suggests semantic pattern detection contributes to raw-leakage containment, while policy, fanout, and safe-view mechanisms still matter. It does not prove that semantic detection is unnecessary.
